\section{Actual stabilization of growing canonical sequences}

Bierstone--Milman\cite{bierstone-milman08-functoriality}, Section 7, and
W{\l}odarczyk\cite{wlodarczyk05-hironaka} use local finiteness of the recorded
centres to obtain eventual local stabilization. This section proves that
bridge for genuine growing canonical sequences. It does not assume, or
assert as a consequence of local finiteness alone, that such a globally
resolving sequence has already been constructed on every noncompact manifold.

\begin{definition}[The actual next centre and its proper image]
\label{hir:growing-sequence-actual-centres}
\uses{hir:centre-function-seq,bcalc:restrict-empty-tower}
\lean{Hironaka.HironakaLocallyFiniteResolutionSequence.stepCentre,Hironaka.HironakaLocallyFiniteResolutionSequence.centreImage,Hironaka.HironakaLocallyFiniteResolutionSequence.exists_step_map}
\leanok
Let $s$ be a growing sequence $X=X_0\leftarrow X_1\leftarrow\cdots$ of
actual recorded blowup towers $T_i$. Choose the actual centre $C_i\subset X_i$
and blowup $\pi_i:X_{i+1}\to X_i$ from its one-step extension, so that
$T_{i+1}=T_i\mathbin{\cdot}(C_i,\pi_i)$. Put $B_i=T_i(C_i)\subset X$.
\end{definition}
\begin{proof}
The growing sequence supplies an actual one-step extension at every index.
Choice selects its centre, and the same witness provides the map, its blowup
property, and the exact equality of recorded towers. The definition of $B_i$
is the image of that actual centre, not an exceptional-set surrogate.
\end{proof}

\begin{lemma}[Counting actual centre intersections]
\label{hir:growing-sequence-count-step}
\uses{hir:growing-sequence-actual-centres,hir:centre-function-seq}
\lean{Hironaka.HironakaLocallyFiniteResolutionSequence.countCentresMeeting_step,Hironaka.HironakaLocallyFiniteResolutionSequence.monotone_countCentresMeeting}
\leanok
For $K\subset X$, write $q_i(K)$ for the number of recorded centres meeting
the inverse image of $K$ in $T_i$. Then
\[
 q_{i+1}(K)=q_i(K)+\mathbf 1_{B_i\cap K\ne\varnothing}.
\]
In particular $q_i(K)$ is nondecreasing.
\end{lemma}
\begin{proof}
Rewrite $T_{i+1}$ using the actual one-step witness. The recursive definition
of the count adds one precisely when $C_i\cap T_i^{-1}(K)$ is nonempty.
This is equivalent to nonemptiness of $B_i\cap K$, by taking the image of a
witness in one direction and its actual preimage witness in the other.
Each increment is either zero or one, so induction on indices gives monotonicity.
\end{proof}

\begin{theorem}[Compact sets eventually meet no later centre]
\label{hir:growing-canonical-sequence-compact-tail}
\uses{hir:growing-sequence-count-step,hir:centre-function-seq}
\lean{Hironaka.HironakaLocallyFiniteResolutionSequence.exists_no_later_centres_meet_compact,Hironaka.HironakaLocallyFiniteResolutionSequence.finite_centres_meeting_compact}
\leanok
Let $R$ be a resolution invariant with per-sequence local finiteness, and
let every $T_i$ be canonical for an actual marked ideal in its domain. For
every compact $K\subset X$, there is $N$ such that $B_i\cap K$ is empty
for all $i\ge N$. In particular only finitely many indices have centres
whose images meet $K$.
\end{theorem}
\begin{proof}
Per-sequence local finiteness bounds all $q_i(K)$ by an integer $m$. Their
range is a nonempty finite subset of $\{0,\ldots,m\}$, so its maximum is
attained at some $N$. Monotonicity and maximality force $q_i(K)=q_N(K)$
for every $i\ge N$. The count identity then forces every later increment
to be zero, which says exactly $B_i\cap K=\varnothing$. The meeting indices
are contained in $\{0,\ldots,N-1\}$, so they form a finite set.
\end{proof}

\begin{theorem}[A locally finite family of closed proper centre images]
\label{hir:growing-canonical-sequence-closed-images}
\uses{hir:growing-sequence-actual-centres,hir:growing-canonical-sequence-compact-tail}
\lean{Hironaka.HironakaLocallyFiniteResolutionSequence.isClosed_centreImage,Hironaka.HironakaLocallyFiniteResolutionSequence.locallyFinite_centreImages}
\leanok
Every $B_i$ is closed, and the indexed family $(B_i)_{i\in\mathbb N}$ is
locally finite on the base manifold $X$.
\end{theorem}
\begin{proof}
The actual centre is a closed smooth submanifold by the blowup property.
The actual tower projection is proper, hence closed, so $B_i$ is closed.
Every point of a complex manifold has a compact neighbourhood $K$. Only
finitely many $B_i$ meet that neighbourhood, by the preceding theorem.
This proves local finiteness of the indexed family, including repeated or
empty images; finiteness counts actual indices rather than distinct sets.
\end{proof}

\begin{theorem}[An open neighbourhood misses the whole actual tail]
\label{hir:growing-canonical-sequence-open-tail}
\uses{hir:growing-canonical-sequence-closed-images,hir:growing-canonical-sequence-compact-tail}
\lean{Hironaka.HironakaLocallyFiniteResolutionSequence.exists_open_no_later_centres}
\leanok
For every compact $K\subset X$ there are $N$ and an open $U\supset K$
such that $C_i\cap T_i^{-1}(U)=\varnothing$ for every $i\ge N$.
\end{theorem}
\begin{proof}
Choose $N$ so all $B_i$ with $i\ge N$ miss $K$. This tail family remains
locally finite. Since every member is closed, its union is closed: near any
point only finitely many members need to be removed. Let $U$ be the complement
of that union. It is open and contains $K$. Any point of $C_i\cap T_i^{-1}(U)$
would have image in both $U$ and the omitted tail union, a contradiction.
\end{proof}

\begin{theorem}[The later actual blowups are isomorphisms over that neighbourhood]
\label{hir:growing-canonical-sequence-isomorphism-tail}
\uses{hir:growing-canonical-sequence-open-tail,bcalc:restrict-empty-tower}
\lean{Hironaka.HironakaLocallyFiniteResolutionSequence.isIsomorphismOver_step_of_disjoint,Hironaka.HironakaLocallyFiniteResolutionSequence.exists_open_isomorphism_tail}
\leanok
For the same $N,U$, every actual later map $\pi_i:X_{i+1}\to X_i$ is an
isomorphism over $T_i^{-1}(U)$. Thus its restriction over $U$ has empty
centre, and every finite later composite stabilizes there up to actual
isomorphism.
\end{theorem}
\begin{proof}
The actual blowup is an isomorphism off $C_i$. Since $T_i^{-1}(U)$ misses
$C_i$, restriction of that actual isomorphism to $T_i^{-1}(U)$ remains
bijective and avoids the exceptional locus. The one-step tower witness
retains the actual map and blowup proof. Successive restricted maps therefore
compose to isomorphisms over $U$; the restricted-centre trace is empty at
all those steps.
\end{proof}


\begin{definition}[The actual finite tail tower]
\label{hir:growing-sequence-actual-tail-tower}
\uses{hir:growing-sequence-actual-centres}
\lean{Hironaka.HironakaLocallyFiniteResolutionSequence.tailTower,Hironaka.HironakaLocallyFiniteResolutionSequence.append_tailTower}
\leanok
For every $N,j$, the actual steps of the growing sequence construct a finite
tower $T_{N,j}:X_N\leftarrow\cdots\leftarrow X_{N+j}$ with
$T_N\mathbin{\cdot}T_{N,j}=T_{N+j}$ as an equality of recorded towers.
\end{definition}
\begin{proof}
For $j=0$ take the empty tower on $X_N$. For $j+1$ append to $T_{N,j}$ the
actual chosen next blowup at index $N+j$. By the one-step witness this is a
genuine blowup and has last manifold $X_{N+j+1}$. Associativity of appending
and the induction equality give
$T_N\mathbin{\cdot}T_{N,j+1}=T_{N+j}\mathbin{\cdot}(C_{N+j},\pi_{N+j})
=T_{N+j+1}$. Thus the construction retains the actual intermediate spaces,
centres and maps, not just their composite.
\end{proof}

\begin{theorem}[The entire actual restricted tail is empty]
\label{hir:growing-sequence-empty-restricted-tails}
\uses{hir:growing-sequence-actual-tail-tower,hir:growing-canonical-sequence-open-tail,bcalc:restrict-empty-tower}
\lean{Hironaka.HironakaLocallyFiniteResolutionSequence.centresAvoid_tailTower,Hironaka.HironakaLocallyFiniteResolutionSequence.exists_open_empty_restricted_tails}
\leanok
The compact-tail neighbourhood $U$ may be chosen so that, for every finite
$j$, the restriction of the actual $T_{N,j}$ to $T_N^{-1}(U)$ has only
empty centres.
\end{theorem}
\begin{proof}
Induct through the actual steps of $T_{N,j}$. The empty tail has no centres.
At the next step, the append equality of recorded towers identifies
$T_N\circ T_{N,j}$ with the original projection $T_{N+j}$. Consequently
the inverse image of $T_N^{-1}(U)$ at that step equals $T_{N+j}^{-1}(U)$.
The actual next centre misses this set by the open-tail theorem. This proves
that all actual tail centres avoid the starting open preimage. The genuine
restricted-tower calculus identifies such avoidance exactly with having
empty centre at every recorded step of the restriction.
\end{proof}

\begin{remark}[The remaining noncompact existence obligation]
To produce a globally resolving sequence one must still glue genuine local
canonical towers and prove a scheduling rule that eventually resolves every
point. The local driver shrinks open bases around compact sets. Its independently
chosen output towers are not already compatible on overlaps. Moreover a
global maximum of the invariant need not be attained on a noncompact
cosupport, so the compact dependent-choice driver cannot simply be reused.
Neither a caller-supplied compatible-tower certificate nor the local finiteness
of a given sequence discharges this remaining obligation.
\end{remark}
